\chapter{Conclusions and future work}
\label{chap:conclusions}

\section{Conclusions}

The general objective was to develop a web application that automates the collection, processing, analysis and visualization of data from YouTube videos. OpenTube Insights runs every service with Docker Compose and covers the complete path from a research question written in natural language to a chart and an exportable dataset. Its contribution is not the collection itself, which existing tools already offer, but that every step of the path leaves a record that another researcher can inspect. The conclusions for each specific objective follow.

\paragraph{Design of the architecture (objective~\ref{obj:design}).} The review of the YouTube Data API showed that search results change over time and that API data may be stored for at most 30 days, so a dataset collected from the API cannot be reproduced identically. The architecture therefore treats the documented collection procedure as the reproducible unit: saved queries, effective parameters, request logs and processing reports are kept permanently and contain no YouTube data, while the YouTube data itself is deleted after 30 days. This separation shaped the database schema and every later component.

\paragraph{Data collection with a natural language interface (objective~\ref{obj:implement_collection}).} Restricting the LLM to a fixed set of typed fields, and building the API parameters in the application, made the translation predictable: the model cannot produce an invalid parameter name, and a single validator applied to drafts, saved queries and run starts keeps invalid values from consuming quota. Keeping the researcher in the loop, with an editable draft and a record of its provenance, preserves the researcher's authority over the criteria. On the collection side, a single logged entry point for all API calls made quota accounting and auditing a property of the architecture rather than a discipline of each caller.

\paragraph{Processing, analysis and visualization (objective~\ref{obj:implement_processing_analysis_and_visualization}).} Comment text is normalized by versioned, deterministic rules, and each processing run reports its inputs, outputs and errors with counts that outlive the purge. Sentiment analysis was built against a contract that records the model name and version with every result, so results of different models are never mixed in the run detail, the chart or the exports. % TODO(US-025 part B): name the model and summarize what its integration showed.
The presentation layer shows each figure together with its context: the model behind a label, the period a count covers and the filters that produced it.

\paragraph{Testing (objective~\ref{obj:test}).} The backend and frontend suites, with 365 and 104 tests, cover each component without calling external services, and an integration test runs the complete workflow on a synthetic corpus. Requiring every test to fail when its behavior is broken on purpose found defects that review alone had missed, such as a pagination loop that could exhaust the daily quota. % TODO(US-039): summarize the validation on the two real topics.

Finally, the development process followed the same principle as the application: with an AI coding agent implementing the stories, the approved plans, review logs and deferred-work records keep every decision attributable and open to inspection.

\section{Limitations}

\begin{itemize}
    \item \textbf{Reproducibility of the data.} The datasets themselves cannot be reproduced, only the procedure, because YouTube search results change over time and the data must be deleted after 30 days.

    \item \textbf{Validation of form, not existence.} Parameters are checked for their format but not against external lists; a nonexistent country code with a valid form is accepted.

    \item \textbf{Local, single-user deployment.} The API has no authentication or rate limiting and attributes requests to a configured local owner, which is adequate only while the application listens on the local machine.

    \item \textbf{Operator duties.} Collection runs and the 30-day purge are executed by management commands that the operator must run or schedule, and the purge must not run during a collection.

    \item \textbf{Interface testing.} The interface logic is tested, but the components are not rendered by automated tests. % TODO(US-037): update once the smoke tests exist.
\end{itemize}

\section{Future work}

\begin{itemize}
    \item Add authentication, per-user quotas and rate limiting, so the application can be deployed as a shared service for a research group.

    \item Schedule collection runs and the purge within the application, with a lock that prevents them from overlapping.

    \item Compare several sentiment models on the same comments, which the data model already supports because every result records its model and version.

    \item Extend the analysis to the channel level, using the channel metadata already collected, and to the evolution of engagement across runs, using the statistics snapshots each run records.

    \item Resolve the open decisions recorded during development, such as preserving emoji composed with the zero-width joiner during normalization.
\end{itemize}
